\documentclass[11pt]{article}
\usepackage[margin=1.1in]{geometry}
\usepackage{amsmath,amssymb,amsthm}
\usepackage[hidelinks]{hyperref}
\usepackage{microtype}
\newtheorem{theorem}{Theorem}[section]
\newtheorem{lemma}[theorem]{Lemma}
\newtheorem{proposition}[theorem]{Proposition}
\newtheorem{corollary}[theorem]{Corollary}
\theoremstyle{definition}
\newtheorem{remark}[theorem]{Remark}
\newtheorem{conjecture}[theorem]{Conjecture}
\newcommand{\N}{\mathbb{N}}
\newcommand{\Lam}{\Lambda}
\DeclareMathOperator{\vp}{v}
\DeclareMathOperator{\om}{\omega}

\title{Explicit upper bounds for the Erd\H{o}s--Sur\'anyi function $g(n)$}
\author{Haoyu Chen\thanks{%
\textbf{Disclosure of the production pipeline.} Every mathematical statement in this note was produced by
automated systems inside an autonomous research pipeline: the proofs of Theorems~\ref{thm:log} and \ref{thm:sqrt} were found in two independent runs
of OpenAI GPT-5.6 at ``Pro'' effort through the web interface (26 and 90 minutes of reasoning respectively,
on 2026-09-03) and the proofs of Theorems~\ref{thm:long} and \ref{thm:ksplit} by Qwen3.8-Max (Alibaba) through its web
interface in two runs on 2026-09-03/04, all working from briefs written by an orchestrating agent (Claude Fable 5.1, Anthropic) that
selected the problem, pinned the statement to Erd\H{o}s's 1992 text, verified the lower-bound construction
by exact computation, re-derived every step of both proofs, commissioned independent referee runs
(Claude Opus 5 agents in fresh contexts, which also brute-forced the lemmas on hundreds of thousands of
instances), and drafted this document. The four constructions were implemented and exercised on more than
seven million random instances, and every numerical inequality in
the explicit parameter choices was checked by machine (scripts in the accompanying repository). No human
supplied a definition, a lemma, a proof idea, or a correction. The problem itself remains open: the
conjectured bound $g(n)\le 2n$ is not proved here.}}
\date{September 4, 2026 (v3; v1: September 3, v2: September 4)}

\begin{document}
\maketitle

\begin{abstract}
For $n\ge 1$ let $g(n)$ be the least integer $g$ such that for every set $\{a_1<\dots<a_n\}$ of integers
greater than $1$ and every $x\ge 0$, some $g$ of the integers $x+1,\dots,x+a_n$ have product divisible by
$a_1\cdots a_n$. Erd\H{o}s and Sur\'anyi (1959) proved $g(3)=4$ and $g(n)\ge(2-o(1))n$, and Erd\H{o}s asked
(offering \$100) whether $g(n)\le(2+o(1))n$, or even $g(n)\le 2n$ (Erd\H{o}s Problem \#708). To our
knowledge no upper bound depending on $n$ alone has been published. We prove
\[
g(n)\;\le\; n\bigl(\lceil\log_2 2n\rceil+\lceil\log_2\lceil\log_2 2n\rceil\rceil+2\bigr)\quad(n\ge 2),
\qquad
g(n)\;\le\;\Bigl\lceil 48\,n\sqrt{\tfrac{81+\ln n}{\ln(81+\ln n)}}\Bigr\rceil ,
\]
and $g(n)\le(2+o(1))\,n\sqrt{\ln n/\ln\ln n}$. We also show that the conjectured bound holds for long
intervals: if $a_n\ge 8n^3$ then $2n$ integers always suffice, so Erd\H{o}s's question is reduced to the
range $a_n<8n^3$; more generally, for every integer $k\ge 2$, $kn$ integers suffice whenever $a_n\le kn$ or
$a_n^{k-1}\ge(kn)^{k+1}$; hence for every fixed $\delta>0$ a bound linear in $n$ (with a constant depending on
$\delta$) holds whenever $a_n\ge n^{1+\delta}$, e.g.\ $12n$ integers suffice whenever $a_n\ge n^2$. All proofs are elementary and constructive: prime-power atoms of the $a_i$ that are large
relative to $a_n$ are placed one per interval element by a per-prime greedy rule, the small atoms are packed
into bins, and a counting argument bounds the number of bins. We also record $g(4)\ge 5$ and $g(5)\ge 6$ by
exact computation.
\end{abstract}

\section{Introduction}

Gallai observed that among any $a_2$ consecutive integers one can always find two whose product is a
multiple of $a_1a_2$ (for $1<a_1<a_2$), but that there are $a_1<a_2<a_3$ and a run of $a_3$ consecutive
integers among which no three have product divisible by $a_1a_2a_3$. Erd\H{o}s and Sur\'anyi
\cite{ES59} studied the general situation. Following Erd\H{o}s's formulation \cite[\S1]{Er92}, let $g(n)$ be
the smallest integer such that, whenever $1<a_1<\dots<a_n$ are integers and $x\ge 0$, there are integers
$x<u_1<\dots<u_{g(n)}\le x+a_n$ with $u_1\cdots u_{g(n)}\equiv 0 \pmod{a_1\cdots a_n}$. (Since enlarging $B$ inside the interval preserves divisibility, ``$g(n)$ integers'' may be read as ``at most
$g(n)$ integers''; we use the latter, which also covers the case $a_n<g(n)$.) They proved
$g(3)=4$ and
\begin{equation}\label{eq:ES}
 g(n)\;\ge\;(2-\varepsilon)\,n\qquad(n>n_0(\varepsilon)),
\end{equation}
the lower bound coming from the set of all products $p_ip_j$ of $\ell$ primes satisfying $2p_1^2>p_\ell^2$,
with the interval positioned so that a single integer $u$, divisible exactly once by each $p_i$, is the only
common multiple of every $p_ip_j$ in it. Erd\H{o}s \cite{Er92} asked whether $g(n)<(2+\varepsilon)n$ for
large $n$, ``or perhaps even $g(n)\le 2n$'', and offered \$100 for a proof or disproof; the question is
Problem \#708 on the Erd\H{o}s problems site \cite{EP708}. Neither \cite{Er92} nor \cite{EP708} records
an upper bound for $g(n)$, and we are not aware of one in the literature (the original paper \cite{ES59},
in Hungarian, was not accessible to us). Note that a bound in terms of $n$ alone is not obvious: the trivial
estimate $g\le\sum_i\om(a_i)$ (one interval element per prime factor of each $a_i$, see
Lemma~\ref{lem:large}) depends on $a_n$, which is unbounded.

Throughout, $\vp_p(k)$ is the exponent of the prime $p$ in $k$, $\om(k)$ the number of distinct prime factors
of $k$, and $\Lam(t)=\lceil\log_2 t\rceil$ is the least $r\ge 0$ with $t\le 2^r$.

\begin{theorem}\label{thm:log}
For every $n\ge 2$,
\[
 g(n)\;\le\;F(n):=n\bigl(\Lam(2n)+\Lam(\Lam(2n))+2\bigr)\;=\;n\log_2 n+n\log_2\log_2 n+O(n).
\]
In particular $F(2)=10$, $F(3)=21$, $F(4)=28$, $F(10)=100$, $F(100)=1300$, $F(1000)=17000$.
\end{theorem}

\begin{theorem}\label{thm:sqrt}
For every $n\ge 1$, with $L=81+\ln n$,
\[
 g(n)\;\le\;\Bigl\lceil 48\,n\sqrt{L/\ln L}\,\Bigr\rceil ,
\qquad\text{and}\qquad
 g(n)\;\le\;(2+o(1))\,n\sqrt{\frac{\ln n}{\ln\ln n}}\quad(n\to\infty).
\]
\end{theorem}

Theorem~\ref{thm:log} gives the smaller value for every $n$ below about $10^{104}$; Theorem~\ref{thm:sqrt}
is asymptotically stronger. Neither is close to the conjectured $2n$. For long intervals, however, the
conjecture holds:

\begin{theorem}\label{thm:long}
If $a_n\ge 8n^3$, then for every $x\ge 0$ some $2n$ of the integers $x+1,\dots,x+a_n$ have product divisible
by $a_1\cdots a_n$ (at most $2n$ are needed, and since $a_n\ge 2n$ the set may be padded); in fact
Conjecture~\ref{conj:S} below holds for such $A$. Consequently $g(n)\le 8n^3$ as
well, and the question whether $g(n)\le 2n$ is equivalent to the same question restricted to $a_n<8n^3$.
\end{theorem}

Theorem~\ref{thm:long} is the case $k=2$ of a one-parameter family:

\begin{theorem}\label{thm:ksplit}
Let $k\ge 2$ be an integer and suppose that $a_n\le kn$ or $a_n^{\,k-1}\ge(kn)^{k+1}$. Then for every $x\ge 0$ some (at
most) $kn$ of the integers $x+1,\dots,x+a_n$ have product divisible by $a_1\cdots a_n$. In particular $a_n\ge 9n^2$
gives $3n$ and $a_n\ge 1024^{1/3}n^{5/3}$ gives $4n$; for every $\delta>0$ and every integer $k$ with
$\delta(k-1)>2$, $a_n\ge n^{1+\delta}$ gives $kn$ as soon as $n^{\delta(k-1)-2}\ge k^{k+1}$. Taking for each $n$ the
smallest admissible $k$ (and Theorem~\ref{thm:log} where none is admissible) gives, for every $n\ge 1$: $7n$ integers
suffice whenever $a_n\ge n^3$, $12n$ whenever $a_n\ge n^2$, and $15n$ whenever $a_n\ge n^{7/4}$.
\end{theorem}

Thus for every fixed $\delta>0$ a bound linear in $n$ holds when $a_n\ge n^{1+\delta}$, with a constant depending on
$\delta$; for a fixed constant $C$ the range not covered by Theorem~\ref{thm:ksplit} is $Cn<a_n<(Cn)^{(C+1)/(C-1)}$,
a full power of $n$. The Erd\H{o}s--Sur\'anyi family, with $a_n\asymp n(\ln n)^2$, lies just above the trivial range.
The gap between $n\sqrt{\ln n/\ln\ln n}$ and $2n$ is discussed in Section~\ref{sec:remarks}.

Finally, Section~\ref{sec:linear} reduces a linear bound to a single inequality about additive functions on
intervals (Conjecture~\ref{conj:TH}): if it holds, then $g(n)\le18n$ (Theorem~\ref{thm:cond}). We prove the
inequality in two special cases and show, with explicit certificates, that three natural strengthenings of it are
false because intervals can contain more integers free of small primes than $[1,m]$ does.

\section{The common construction}\label{sec:construction}

Fix $A=\{a_1<\dots<a_n\}\subset\{2,3,\dots\}$, put $m=a_n$, $I=\{x+1,\dots,x+m\}$ and $P=a_1\cdots a_n$.
We must find $B\subseteq I$ with $P\mid\prod_{b\in B}b$ and bound $|B|$. Note that $B$ is a set: an interval
element contributes its valuations once, however many $a_i$ it may be a multiple of.

\begin{lemma}[counting]\label{lem:count}
For every integer $d\le m$: $\;\#\{a\in A: d\mid a\}\;\le\;\lfloor m/d\rfloor\;\le\;\lfloor (x+m)/d\rfloor-\lfloor x/d\rfloor=\#\{b\in I: d\mid b\}$.
\end{lemma}
\begin{proof}
The multiples of $d$ in $\{1,\dots,m\}$ are $d,2d,\dots,\lfloor m/d\rfloor d$, and the $a_i$ are distinct
elements of $\{2,\dots,m\}$. Any $m$ consecutive integers contain
$\lfloor (x+m)/d\rfloor-\lfloor x/d\rfloor\ge\lfloor m/d\rfloor$ multiples of $d$.
\end{proof}

\begin{lemma}[whole interval]\label{lem:whole}
$P$ divides $\prod_{b\in I}b$. Hence if $m\le H$ then $B=I$ has $|B|=m\le H$.
\end{lemma}
\begin{proof}
$P\mid m!$ because the $a_i$ are distinct elements of $\{2,\dots,m\}$, and
$\prod_{b\in I}b=m!\binom{x+m}{m}$.
\end{proof}

\paragraph{Atoms, large and small.}
Fix an integer budget $H\ge 1$ and assume $m>H$; put $T=m/H>1$. For $a\in A$ and $p\mid a$ call
$q_p(a)=p^{\vp_p(a)}$ an \emph{atom} of $a$; the atoms of $a$ are pairwise coprime and multiply to $a$, and
$a$ has $\om(a)$ atoms. An atom is \emph{large} if $q>T$ and \emph{small} if $q\le T$.

\begin{lemma}[large atoms, one element each]\label{lem:large}
Let $\mathcal L$ be the multiset of all large atoms of all $a\in A$ (atoms of different $a$ are distinct
members, so $r=|\mathcal L|$ counts with multiplicity). There is a \emph{set} $U\subseteq I$ with $|U|\le r$
such that for every prime $p$
\[
 \vp_p\Bigl(\prod_{u\in U}u\Bigr)\;\ge\;\sum_{q\in\mathcal L,\;q\text{ a power of }p}\vp_p(q).
\]
More precisely, one can assign to each $q\in\mathcal L$ an element $\beta(q)\in I$ with $q\mid\beta(q)$ such
that atoms of the same prime receive distinct elements; then $U=\beta(\mathcal L)$.
\end{lemma}
\begin{proof}
Fix a prime $p$ and list the large atoms of $p$ as $p^{e_1},\dots,p^{e_s}$ with $e_1\ge\dots\ge e_s$. Each
$a\in A$ has at most one atom of $p$, so the first $j$ atoms come from $j$ distinct elements of $A$, all
divisible by $p^{e_j}$; by Lemma~\ref{lem:count}, $I$ contains at least $j$ multiples of $p^{e_j}$. Assign the
atoms in this order, giving the $j$-th one a multiple of $p^{e_j}$ in $I$ different from the $j-1$ elements
already used for $p$. Elements used for different primes may coincide; this only decreases $|U|$. For a
fixed $p$ the assigned elements are distinct, so their $p$-adic valuations add up to at least $\sum_j e_j$.
\end{proof}

\begin{lemma}[assembly]\label{lem:assembly}
Suppose the small atoms of all $a\in A$ (a multiset; atoms of different $a$ are distinct members) have been
partitioned into \emph{bins}, each bin being a set of small atoms whose product $f$ satisfies $f\le T$ (a bin may
contain atoms of different $a$). Let $s$ be the total number of bins and
$D=r+s$. If $D\le H$, then there is $B\subseteq I$ with $|B|\le D$ and $P\mid\prod_{b\in B}b$.
\end{lemma}
\begin{proof}
Take $U$ from Lemma~\ref{lem:large}, $|U|\le r$. Order the bins $f_1,\dots,f_s$ arbitrarily. Since
$f_j\le T=m/H$ and $H$ is an integer, Lemma~\ref{lem:count} gives at least $\lfloor m/f_j\rfloor\ge H$ multiples of
$f_j$ in $I$.
When bin $j$ is processed, the elements to avoid are those of $U$ and the $j-1$ elements chosen for earlier
bins, at most $r+s-1=D-1\le H-1$ in number; so an unused multiple $v_j$ of $f_j$ exists. Put
$B=U\cup\{v_1,\dots,v_s\}$. By construction $U\cap\{v_1,\dots,v_s\}=\emptyset$ and the $v_j$ are pairwise
distinct, so $|B|\le r+s=D$ and (this is the only place where the valuations of different elements are added)
\[
 \prod_{b\in B}b=\prod_{u\in U}u\cdot\prod_{j}v_j ,
\]
whose $p$-adic valuation is at least (sum over large atoms of $p$) $+\;\vp_p(\prod_j f_j)$, which equals
$\vp_p(P)$ because every atom of every $a$ is either large or lies in exactly one bin.
\end{proof}

Thus everything reduces to producing bins and a budget $H$ with $D=r+s\le H$. Theorems~\ref{thm:log} and
\ref{thm:sqrt} do this in two different ways.

\section{Proof of Theorem~\ref{thm:log}}\label{sec:log}

\paragraph{Next-fit bins.} For each $a\in A$ list its small atoms in increasing order of the prime and pack
them by multiplicative next-fit: keep a current bin, initially empty (product $1$); if multiplying the next
atom keeps the product $\le T$, add it; otherwise close the bin and start a new one containing that atom.
Every small atom is $\le T$, so every bin has product $\le T$.

\begin{lemma}\label{lem:nextfit}
If $d_1,\dots,d_u$ are the bin products of one $a$ in the order created, then $d_jd_{j+1}>T$ for $1\le j<u$.
\end{lemma}
\begin{proof}
Bin $j$ was closed because the next atom $q$ did not fit: $d_jq>T$. That atom is the first element of bin
$j+1$, so $d_{j+1}\ge q$.
\end{proof}

\begin{lemma}[budget]\label{lem:budget}
Let $H,\kappa\ge 1$ be integers with
\begin{equation}\label{eq:B}
 \text{(B1)}\;\; H\ge n(2\kappa+1),\qquad\text{(B2)}\;\; H^{(\kappa+1)n}\le 2^{\kappa H}.
\end{equation}
Then the next-fit construction with threshold $T=m/H$ produces $D=r+s\le H$ demands.
\end{lemma}
\begin{proof}
\emph{Case $m^\kappa\le H^{\kappa+1}$.} Binning never increases the count, so $D\le W:=\sum_{a\in A}\om(a)$.
Since every atom is at least $2$, $2^{\om(a)}\le a$, hence $2^W\le\prod_a a\le m^n$ and
$2^{\kappa W}\le m^{\kappa n}\le H^{(\kappa+1)n}\le 2^{\kappa H}$ by (B2). So $W\le H$.

\emph{Case $m^\kappa>H^{\kappa+1}$.} Then $T^{\kappa+1}=m^{\kappa+1}/H^{\kappa+1}>m$. Fix $a\in A$ with
$\ell$ large atoms and $u$ bins $d_1,\dots,d_u$. The large atoms and the products $d_1d_2,\,d_3d_4,\dots$ of
consecutive bin pairs are pairwise coprime divisors of $a$, each exceeding $T$ (Lemma~\ref{lem:nextfit}).
If $\ell+\lfloor u/2\rfloor\ge\kappa+1$, the product of $\kappa+1$ of them would exceed $T^{\kappa+1}>m\ge a$
while dividing $a$; so $\ell+\lfloor u/2\rfloor\le\kappa$, and since $u\le 2\lfloor u/2\rfloor+1$,
$\ell+u\le 2(\ell+\lfloor u/2\rfloor)+1\le 2\kappa+1$. Summing over $a$, $D\le n(2\kappa+1)\le H$ by (B1).
\end{proof}

\begin{lemma}\label{lem:params}
For $n\ge 2$ let $r=\Lam(2n)$, $s=\Lam(r)$, $C=r+s+2$, $H=nC=F(n)$ and $\kappa=\lfloor(C-1)/2\rfloor$. Then
\textup{(B1)} and \textup{(B2)} hold.
\end{lemma}
\begin{proof}
(B1): $2\kappa+1\le C$. (B2), $n\ge 3$: then $2n\ge 6$, so $r\ge 3$ and $s\ge 2$; by definition $2n\le 2^r$,
i.e.\ $n\le 2^{r-1}$, and $r\le 2^s$; also $s+2\le 2^s$ for $s\ge 2$. Hence $C=r+s+2\le 2^{s+1}$ and
$H=nC\le 2^{r-1}2^{s+1}=2^{C-2}$. Moreover $C\le 2(\kappa+1)$, i.e.\ $(C-2)(\kappa+1)\le\kappa C$, so
$H^{(\kappa+1)n}\le 2^{(C-2)(\kappa+1)n}\le 2^{\kappa Cn}=2^{\kappa H}$. For $n=2$: $r=2$, $s=1$, $C=5$,
$H=10$, $\kappa=2$, and $10^6\le 2^{20}$.
\end{proof}

\begin{proof}[Proof of Theorem~\ref{thm:log}]
Let $H=F(n)$. If $m\le H$ use Lemma~\ref{lem:whole}. Otherwise Lemmas~\ref{lem:budget} and \ref{lem:params}
give $D\le H$, and Lemma~\ref{lem:assembly} gives $B$ with $|B|\le D\le F(n)$.
\end{proof}

\section{Proof of Theorem~\ref{thm:sqrt}}\label{sec:sqrt}

Here the bins are formed by \emph{maximal merging}: start with each small atom (of any $a$) in its own bin
and, as long as two bins have product $\le T$, merge them. Each merge preserves the property ``product
$\le T$'' and decreases the number of bins, so the process terminates, and at termination any two bins have
product $>T$.

\begin{lemma}\label{lem:K}
With $W=\sum_a\om(a)$ and $D=r+s$ as before,
\[
 \text{(i)}\;\; D\le W,\qquad
 \text{(ii)}\;\; D\;\le\;\frac{2n\ln m}{\ln(m/H)}+1 .
\]
\end{lemma}
\begin{proof}
(i) is clear. (ii): any two bins have product $>T$, so at most one bin has product $\le\sqrt T$, and that bin still has
product $\ge 2>1$. Hence, if $s\ge 1$, $P=\prod_{\text{large}}q\cdot\prod_j f_j>T^{\,r}\,T^{(s-1)/2}$ strictly, and
$P\le m^n$; thus
$r+(s-1)/2<n\ln m/\ln T$ and $D=r+s<2n\ln m/\ln T-r+1\le 2n\ln m/\ln(m/H)+1$. If $s=0$, $P>T^{\,r}$ gives
$D=r<n\ln m/\ln T$.
\end{proof}

\begin{lemma}\label{lem:primesum}
For real $X>1$, $\;\sum_{p\le X}1/p\le e\ln(1+\ln X)$ (sum over primes).
\end{lemma}
\begin{proof}
Put $\sigma=1+1/\ln X$. For $p\le X$, $p^{\sigma-1}\le X^{1/\ln X}=e$, so $1/p\le e\,p^{-\sigma}$. Using
$-\ln(1-t)\ge t$ and the finite Euler product,
$\sum_{p\le X}p^{-\sigma}\le\ln\prod_{p\le X}(1-p^{-\sigma})^{-1}\le\ln\zeta(\sigma)$, and
$\zeta(\sigma)\le 1+\int_1^\infty t^{-\sigma}dt=1+1/(\sigma-1)=1+\ln X$.
\end{proof}

\begin{lemma}[average number of prime factors]\label{lem:omega}
For every real $t>1$,
\[
 \frac1n\sum_{i=1}^n\om(a_i)\;\le\;\frac{\ln(m/n)+(t-1)\,e\ln(1+\ln m)}{\ln t}.
\]
\end{lemma}
\begin{proof}
For every positive integer $k$, $t^{\om(k)}=\prod_{p\mid k}\bigl(1+(t-1)\bigr)=\sum_{d\mid k,\ d\text{ squarefree}}(t-1)^{\om(d)}$.
All terms are nonnegative because $t>1$. Summing over $k\le m$,
\[
 \sum_{k\le m}t^{\om(k)}=\sum_{\substack{d\le m\\ d\text{ sqfree}}}(t-1)^{\om(d)}\Bigl\lfloor\frac md\Bigr\rfloor
 \le m\prod_{p\le m}\Bigl(1+\frac{t-1}p\Bigr)\le m\exp\Bigl((t-1)\sum_{p\le m}\frac1p\Bigr)
 \le m\,e^{(t-1)e\ln(1+\ln m)}
\]
by Lemma~\ref{lem:primesum}. The $a_i$ are distinct integers in $[2,m]$, so
$\sum_i t^{\om(a_i)}\le\sum_{k\le m}t^{\om(k)}$, and by the AM--GM inequality
$t^{\frac1n\sum_i\om(a_i)}\le\frac1n\sum_i t^{\om(a_i)}$. Taking logarithms gives the claim (note $m\ge n+1$, so $\ln(m/n)>0$).
\end{proof}

\begin{lemma}[explicit parameters]\label{lem:explicit}
Let $L=81+\ln n$, $\ell=\ln L$, $R=\sqrt{L/\ell}$ and $H=\lceil 48nR\rceil$. If $m>H$ then $D<H$.
\end{lemma}
\begin{proof}
Elementary facts: $\ell>4$ (as $L\ge 81>e^4$); $\ell\le L/4$ (the function $y/4-\ln y$ is increasing and
positive for $y\ge 81$); $\ln\ell\le\ell/2$ (as $y/2-\ln y>0$ is increasing for $y\ge 4$); hence
$R>2$. Put $y=\ln m$, $z=\ln H$, $d=y-z>0$, $c=z-\ln n$. Since $48nR>1$, $H\le 96nR$, so
$c\le\ln 96+\ln R<7+\ell/2$ (using $\ln R=(\ell-\ln\ell)/2\le\ell/2$) and $z<\ln n+7+\ell/2<2L$. Put
$\Delta=4\sqrt{L\ell}=4R\ell\le 2L$ (as $4\ell\le L$).

\emph{Case $d\le\Delta$.} Then $y\le 4L$. Apply Lemma~\ref{lem:omega} with $t=R$: $\ln R=(\ell-\ln\ell)/2\ge\ell/4$
and $e\ln(1+y)\le e\ln(5L)<3(\ell+3)<6\ell$. Since $\ln(m/n)=d+c$,
\[
 \frac Wn\le\frac{d+c+(R-1)\cdot 6\ell}{\ell/4}<\frac{4R\ell+7+\ell/2+6R\ell}{\ell/4}=40R+\frac{28}{\ell}+2<45R,
\]
so by Lemma~\ref{lem:K}(i), $D\le W<45nR<H$.

\emph{Case $d>\Delta$.} By Lemma~\ref{lem:K}(ii),
$D\le 2n\,\frac{z+d}{d}+1<2n+\frac{2n\cdot 2L}{4\sqrt{L\ell}}+1=2n+nR+1<3nR<H$, using $R>2$.
\end{proof}

\begin{proof}[Proof of Theorem~\ref{thm:sqrt}, explicit part]
Let $H=\lceil 48nR\rceil$. If $m\le H$ use Lemma~\ref{lem:whole}; otherwise Lemma~\ref{lem:explicit} gives
$D<H$ and Lemma~\ref{lem:assembly} gives $|B|\le D<H$.
\end{proof}

\begin{proof}[Proof of Theorem~\ref{thm:sqrt}, asymptotic part]
Fix $\varepsilon>0$, put $Q_n=\sqrt{\ln n/\ln\ln n}$ and $H=\lceil(2+\varepsilon)nQ_n\rceil$; we show $D<H$
for all $m>H$ once $n$ is large, which suffices as before. With $y=\ln m$, $z=\ln H$, $d=y-z$, $c=z-\ln n$ we
have $z=\ln n+O(\ln\ln n)$, hence $\sqrt{z/\ln z}=(1+o(1))Q_n$ and $c=O(\ln z)$. Put $\Delta=\sqrt{z\ln z}$.
If $d>\Delta$, Lemma~\ref{lem:K}(ii) gives $D/n\le 2+2z/d+1/n\le(2+o(1))\sqrt{z/\ln z}$. If $d\le\Delta$, then
$y=(1+o(1))z$; apply Lemma~\ref{lem:omega} with $t=\sqrt z/(\ln z)^2$, so $\ln t=(\tfrac12+o(1))\ln z$:
the three terms of the numerator divided by $\ln t$ are at most $(2+o(1))\sqrt{z/\ln z}$, $O(1)$ and
$O(\sqrt z/(\ln z)^2)$ respectively, so $D/n\le W/n\le(2+o(1))\sqrt{z/\ln z}$. In both cases
$D\le(2+\varepsilon/2)nQ_n<H$ for all $n\ge n_0(\varepsilon)$, so $g(n)\le\lceil(2+\varepsilon)nQ_n\rceil$ for
$n\ge n_0(\varepsilon)$; as $\varepsilon>0$ was arbitrary, $\limsup_n g(n)/(nQ_n)\le 2$.
\end{proof}

\begin{remark}
No attempt was made to optimise the constant $48$; the same argument with the case threshold and the
prime-sum estimate tightened (e.g.\ an explicit Mertens bound in place of Lemma~\ref{lem:primesum}) gives
a single-digit constant. Theorems~\ref{thm:log} and \ref{thm:sqrt} are proved by the same construction and
differ only in how the number of bins is counted; Theorem~\ref{thm:long} uses the same two ingredients
(per-prime placement of large atoms, private multiples for small parts) with the budget $2n$ made possible by
the length of the interval.
\end{remark}

\section{Long intervals: proof of Theorem~\ref{thm:long}}\label{sec:long}

Assume $m=a_n\ge 8n^3$ and put $C=m/(2n)$. Then $C\ge m^{2/3}$ (equivalent to $m^{1/3}\ge 2n$),
$C\ge 4n^2\ge 2n$, and, since $m\ge 8n^3\ge 4n^2$, also $C^2=m^2/(4n^2)\ge m$. In the language of
Section~\ref{sec:construction} we take the integer budget $H=2n$ and threshold $T=C=m/H$ (note $m>H$).

\begin{lemma}[two-bin packing]\label{lem:pack}
Let $C\ge1$ be real and let $q_1,\dots,q_r>1$ be integers with $q_i\le C$ and $P=q_1\cdots q_r\le C^{3/2}$.
Then some sub-product $s$ of the $q_i$ (over a subset, possibly empty or everything) satisfies
$P/C\le s\le C$; consequently the $q_i$ can be partitioned into two classes with products $\le C$.
\end{lemma}
\begin{proof}
If $P\le C$ take $s=P$. Otherwise $1<P/C\le\sqrt C$. If some $q_i>\sqrt C$, take $s=q_i$: then $s\le C$ and
$s>\sqrt C\ge P/C$. If all $q_i\le\sqrt C$, multiply them in any order and stop at the first partial product
$s\ge P/C$ (this happens at the latest at $s=P$); the previous partial product, possibly the empty product
$1<P/C$, was $<P/C$, so $s<(P/C)\sqrt C=P/\sqrt C\le C$. In every case the
complementary product is $P/s\le C$.
\end{proof}

\begin{lemma}[split]\label{lem:split}
Every $a\in A$ has at most one atom exceeding $C$, and admits a factorisation $a=d_1d_2$ into coprime parts
such that either $d_1,d_2\le C$, or $d_1$ is an atom $>C$ and $d_2<2n\le C$.
\end{lemma}
\begin{proof}
Two atoms $>C$ are coprime, so they would give $a>C^2\ge m$, impossible. If $q=p^e>C$ is an atom, put
$d_1=q$, $d_2=a/q<m/C=2n$. Otherwise all atoms of $a$ are $\le C$ and their product $a\le m\le C^{3/2}$, so
Lemma~\ref{lem:pack} splits them into two coprime classes with products $\le C$.
\end{proof}

\begin{proof}[Proof of Theorem~\ref{thm:long}]
Split every $a\in A$ by Lemma~\ref{lem:split}; parts equal to $1$ impose no condition and may be assigned
to any element, so we discard them. This gives at most $2n$ \emph{demands}, each either a \emph{large} atom
$>C$ or a \emph{small} integer $\le C$. (Lemma~\ref{lem:assembly} with $H=2n$ already gives $|B|\le 2n$; we
repeat the argument because Conjecture~\ref{conj:S} asks for the explicit assignment.) Assign the large
demands as in Lemma~\ref{lem:large}: for each prime $p$, the large atoms of $p$ (one per $a$) receive
distinct multiples of the required powers of $p$; different primes may share elements. Let $U$ be the set of
elements used, $|U|\le$ number of large demands. Now assign the small demands one at a time: a small demand
$d\le C=m/(2n)$ has at least $\lfloor m/d\rfloor\ge 2n$ multiples in $I$ by Lemma~\ref{lem:count} ($2n$ is
an integer), while the forbidden elements ($U$ together with the elements already given to earlier small
demands) number at most $2n-1$, because there are at most $2n$ demands in all and the current one is not yet
placed; so $d$ receives a multiple of its own, outside $U$ and different from all other small
representatives. Let $B$ be the set of all elements used, $|B|\le 2n$. At an element of $U$ the demands
present are prime powers of distinct primes, each dividing the element; at every other element of $B$
exactly one demand $d$ is present and $d$ divides the element. Hence the product of the demands at any
$b\in B$ divides $b$, and since the demands of each $a$ multiply to $a$, $P\mid\prod_{b\in B}b$. This is
exactly an assignment as in Conjecture~\ref{conj:S}. Finally, if $a_n<8n^3$ then Lemma~\ref{lem:whole}
gives a solution with at most $a_n<8n^3$ elements, so $g(n)\le 8n^3$.
\end{proof}

\section{The $k$-split lemma: proof of Theorem~\ref{thm:ksplit}}\label{sec:ksplit}

\begin{lemma}[$k$-split]\label{lem:ksplit}
Let $k\ge 1$ and let $H\le m$ be positive integers. Let $Q$ be a finite multiset of integers $>1$ with product $A$
such that $A^2H^{k+1}\le m^{k+1}$. Then $Q$ can be partitioned into at most $k$ pieces (none if $Q=\emptyset$), each
of which is either a single element of $Q$ or a \emph{bin}: a nonempty $F\subseteq Q$ whose product $f$ satisfies
$fH\le m$.
\end{lemma}
\begin{proof}
Induction on $k$; if $Q=\emptyset$ there is nothing to do, so let $Q\ne\emptyset$. For $k=1$ the hypothesis reads
$A^2H^2\le m^2$, so $AH\le m$ and $Q$ itself is one bin. Let $k\ge 2$.

\emph{Case 1: some $q\in Q$ has $q^2H>m$.} Make $q$ a singleton piece. The product $A'=A/q$ of the rest satisfies
$A'^2H^k=A^2H^{k+1}/(q^2H)<m^{k+1}/m=m^k$, so the induction hypothesis (with $k-1$) splits the rest into at most
$k-1$ pieces.

\emph{Case 2: $q^2H\le m$ for all $q\in Q$.} Build one bin greedily: start with $f=1$ and, while $Q$ is not exhausted
and $f^2H<m$, multiply the next element into $f$. The invariant $fH\le m$ holds at the start ($H\le m$) and is
preserved, because before a step $f^2H<m$ and $q^2H\le m$ give $(fqH)^2<m^2$. The first element is always absorbed
unless $H=m$, in which case the hypothesis forces $A=1$, i.e.\ $Q=\emptyset$; so the bin is nonempty. If the process
exhausts $Q$, all of $Q$ is one bin. Otherwise it stops with $f^2H\ge m$; then the product $A'=A/f$ of the remaining
elements satisfies $A'^2H^k=A^2H^{k+1}/(f^2H)\le m^k$, and the induction hypothesis splits them into at most $k-1$
further pieces.
\end{proof}

\begin{proof}[Proof of Theorem~\ref{thm:ksplit}]
Put $H=kn$. If $m\le H$, Lemma~\ref{lem:whole} gives $B=I$ with $|B|=m\le kn$. Otherwise $m>H$ and
$m^{k-1}\ge H^{k+1}$; for every $a\in A$, $a^2H^{k+1}\le m^2\cdot m^{k-1}=m^{k+1}$. Apply Lemma~\ref{lem:ksplit} to
the atoms of each $a\in A$ (pairwise coprime, product $a$, all $\le m$): each $a$ splits into at most $k$ pieces, hence
at most $kn=H$ pieces in all. Call a singleton piece $q$ \emph{large} if $qH>m$; every other piece (a bin, or a
singleton with $qH\le m$) is a bin $f$ with $fH\le m$ in the sense of Section~\ref{sec:construction} with threshold
$T=m/H$. The large pieces are atoms of distinct elements of $A$ for each fixed prime (an element of $A$ has one atom
per prime), so Lemma~\ref{lem:large} applies to them, and Lemma~\ref{lem:assembly} (with $r$ = number of large pieces,
$s$ = number of bins, $r+s\le H$) gives $B\subseteq I$ with $|B|\le H=kn$ and $P\mid\prod_{b\in B}b$; note that the
$p$-adic valuation of a bin comes from the single $p$-atom of the element of $A$ it was cut from.
For the corollaries: $k=2$ is $m\ge 8n^3$; $k=3$ is $m^2\ge 81n^4$; $k=4$ is $m^3\ge 1024n^5$; for $\delta>0$ and
$\delta(k-1)>2$, $m\ge n^{1+\delta}$ gives $m^{k-1}\ge n^{(1+\delta)(k-1)}\ge k^{k+1}n^{k+1}$ as soon as
$n^{\delta(k-1)-2}\ge k^{k+1}$. The last sentence of the theorem is a finite computation: for each exponent $e$ and each
$n$ take the least $k\ge 2$ with $e(k-1)\ln n\ge(k+1)\ln(kn)$, or $F(n)/n$ from Theorem~\ref{thm:log} if there is
none; for $e=3,2,7/4$ some $k\le 10$ is admissible for all $n\ge 5$, $n\ge 38$, $n\ge 207$ respectively, and
the maximum over $n$ is attained at $n=4,36,165$ with the values $7,12,15$; the script is in the repository.
\end{proof}

\begin{remark}
Lemma~\ref{lem:ksplit} is the sharp form of the same atom-binning that underlies Theorem~\ref{thm:log} (which needs no
hypothesis on $a_n$); what is new here is the exact threshold $\sqrt{m/H}$ and the explicit family of thresholds
$a_n^{k-1}\ge(kn)^{k+1}$. The threshold is essentially sharp for the lemma: $k+1$ atoms each with $q^2H>m$ cannot share
bins and violate the hypothesis. With $k=2$ the lemma is Lemma~\ref{lem:pack} in a different normalisation, so
Section~\ref{sec:long} is the case $k=2$; it was kept because it also exhibits the explicit assignment required by
Conjecture~\ref{conj:S}. Choosing $k$ as a function of $m/n$ does not improve the global bounds of
Theorems~\ref{thm:log} and \ref{thm:sqrt}: for $m=n\,e^{\sqrt{\ln n\ln\ln n}}$ both routes give about
$n\sqrt{\ln n/\ln\ln n}$.
\end{remark}

\section{Small values, obstructions and remarks}\label{sec:remarks}

\begin{proposition}\label{prop:small}
$g(4)\ge 5$ and $g(5)\ge 6$.
\end{proposition}
\begin{proof}
Let $A=\{10,12,15,20\}$ and $x=50$, so $I=\{51,\dots,70\}$ and $P=2^5\cdot3^2\cdot5^3$. No element of $I$ is
divisible by $25$, so a solution contains at least three of the multiples $55,60,65,70$ of $5$; any set $S$ of
multiples of $5$ in $I$ has $\vp_2(\prod S)\le 3$ and $\vp_3(\prod S)\le 1$. Hence a further element with
$\vp_2\ge 2$ and $\vp_3\ge 1$, i.e.\ a multiple of $12$, is needed; the only multiple of $12$ in $I$ is $60$,
and using it still leaves $\vp_2\le 3<5$. So at least five elements are needed, and $\{55,63,64,65,70\}$
works. For $A=\{5,10,12,15,20\}$ and $x=50$ we have $P=2^5\cdot3^2\cdot5^4$, so all four multiples of $5$ are
forced, supplying $\vp_2=3$ and $\vp_3=1$; two more elements are then needed for $2^2$ and $3$ (no element of
$I$ other than $60$ is divisible by $12$), and $\{55,60,63,64,65,70\}$ works. Both minima were also confirmed
by an exhaustive dynamic programme over capped valuation vectors, which further shows that no other $x$
gives a larger minimum for these two sets.
\end{proof}

\begin{remark}[why the obvious approaches fail]
(a) Choosing a multiple of each $a_i$ and multiplying does not work, because the chosen multiples may
coincide and a shared element contributes its valuations once; in the Erd\H{o}s--Sur\'anyi family every
$a_i$ has the same unique multiple $u$. (b) One cannot in general match the $a_i\le m/2$ to
\emph{distinct} multiples inside an interval of length $m$: van Doorn, Li and Tang \cite{vDLT26} showed that an
interval of length $2\max A$ guarantees only $\min(n,\lceil 2\sqrt n\rceil)$ disjoint pairs $(a,\text{multiple of }a)$,
and this is optimal. (c) An induction that removes $a_n$, takes a solution $B'$ for the rest and repairs it
with two new elements fails: for $A'=\{2,4,8,10,12\}$, $I=\{1,\dots,16\}$, $B'=\{8,12,15,16\}$ and $a_6=16$,
the missing factor $2^4$ cannot be supplied by two elements of $I\setminus B'$. Lemma~\ref{lem:assembly}
avoids all three pitfalls by giving every bin at least $H$ candidates and letting different primes share an
element only through the per-prime bookkeeping of Lemma~\ref{lem:large}.
\end{remark}

\begin{remark}[the gap to $2n$]
Both proofs pay one interval element per large atom and per bin. In the balancing region of
Lemma~\ref{lem:K} an $a_i$ may have about $\sqrt{\ln n/\ln\ln n}$ atoms of ``medium'' size, none of which can
be merged, and this is exactly the loss in Theorem~\ref{thm:sqrt}. A linear bound would require medium atoms
of \emph{different} $a_i$ to share representatives. A natural sufficient statement is the following.
\end{remark}

\begin{conjecture}[split assignment]\label{conj:S}
For every $A$ and $I$ as above one can choose for each $a\in A$ a factorisation $a=d_1(a)d_2(a)$ into
coprime parts ($d_2=1$ allowed) and assign each of the at most $2n$ parts to an element $\beta(d)\in I$ so
that for every $b\in I$ and every prime $p$, $\sum_{d:\beta(d)=b}\vp_p(d)\le\vp_p(b)$.
\end{conjecture}

The conjecture implies $g(n)\le 2n$ (take $B=\beta(\text{parts})$), and Theorem~\ref{thm:long} proves it
whenever $a_n\ge 8n^3$. Since the capacity constraint couples
two parts only if they share a prime, the assignment always exists when every $a\in A$ has at most two
distinct prime factors (use the finest splitting and Lemma~\ref{lem:large} prime by prime); a counterexample
would need elements with three or more prime factors competing for composite multiples. An exact search found
the assignment on the Erd\H{o}s--Sur\'anyi instances with $\ell=3,4$, on the instances of
Proposition~\ref{prop:small}, and on more than $10^6$ random instances with $n\le 6$ and $a_n\le 210$;
we found no counterexample.

\section{A conditional linear bound}\label{sec:linear}

Theorem~\ref{thm:long} settles the range $a_n\ge 8n^3$ with the conjectured value $2n$, so a bound $g(n)\le Cn$
with an absolute constant $C$ is a statement about $a_n<8n^3$ only. In this section we reduce such a bound to a single
inequality about additive functions on intervals of integers, we prove the inequality in two special cases, and we
record three natural strengthenings of it that are \emph{false}, with explicit certificates.

Throughout, a \emph{weight} is a family $z=(z_p)_p$ indexed by the primes with $0\le z_p\le 1$ and only finitely many
$z_p\ne0$, and $w_z(k):=\sum_p z_p\vp_p(k)$ is the associated completely additive function. We write $t^+=\max(t,0)$.

\begin{conjecture}[hinge inequality]\label{conj:TH}
For every $m\ge1$, every $x\ge0$ and every weight $z$, with $I=\{x+1,\dots,x+m\}$,
\begin{equation}\label{eq:TH}
\sum_{k=1}^{m}\bigl(w_z(k)-2\bigr)^+\;\le\;\sum_{b\in I}\bigl(w_z(b)-1\bigr)^+ .
\end{equation}
\end{conjecture}

\begin{theorem}\label{thm:cond}
If Conjecture~\ref{conj:TH} holds, then $g(n)\le 18n$ for all $n\ge1$. More precisely, it suffices that
\eqref{eq:TH} holds for $m=a_n$ and every weight supported on the primes dividing $\prod A$.
\end{theorem}

The proof is a linear-programming argument. Fix $A=\{a_1<\dots<a_n\}$ with $m=a_n<8n^3$, an interval $I$ of length
$m$, the set $\mathcal P$ of primes dividing $P=\prod A$, and $R_p:=\sum_{a\in A}\vp_p(a)$ for $p\in\mathcal P$.
Consider the fractional cover problem
\begin{equation}\label{eq:LP}
\tau^*:=\min\Bigl\{\sum_{b\in I}y_b:\ \sum_{b\in I}\vp_p(b)\,y_b\ge R_p\ (p\in\mathcal P),\ 0\le y_b\le1\Bigr\}.
\end{equation}
It is feasible, since $y\equiv1$ works: $\prod A$ divides $m!$, which divides the product of any $m$ consecutive
integers, so $\sum_{b\in I}\vp_p(b)\ge R_p$.

\begin{lemma}[duality]\label{lem:dual}
$\tau^*=\max_{z\ge0}\Bigl[\sum_{a\in A}w_z(a)-\sum_{b\in I}\bigl(w_z(b)-1\bigr)^+\Bigr]$, where the maximum is over
$z=(z_p)_{p\in\mathcal P}$ with $z_p\ge0$, and it is attained at some $z$ with all $z_p\le1$.
\end{lemma}
\begin{proof}
The dual of \eqref{eq:LP} has variables $z_p\ge0$ (one per covering constraint) and $u_b\ge0$ (one per upper bound
$y_b\le1$) and reads $\max\sum_p z_pR_p-\sum_b u_b$ subject to $\sum_p z_p\vp_p(b)-u_b\le1$ for every $b\in I$.
For fixed $z$ the best choice is $u_b=(w_z(b)-1)^+$, and $\sum_p z_pR_p=\sum_{a\in A}w_z(a)$; strong duality holds
because \eqref{eq:LP} is feasible and bounded. For the last claim fix all $z_q$, $q\ne p$, and let $z_p\ge1$ vary:
every $b$ with $p\mid b$ then has $w_z(b)\ge1$, so the right derivative of the objective in $z_p$ equals
$R_p-\sum_{b\in I}\vp_p(b)\le0$. Hence the objective does not increase on $z_p\ge1$.
\end{proof}

\begin{lemma}[rounding]\label{lem:round}
There is a set $B\subseteq I$ with $\prod A\mid\prod B$ and $|B|\le\tau^*+|\mathcal P|$.
\end{lemma}
\begin{proof}
Let $y^*$ be an optimal vertex of the polytope in \eqref{eq:LP}. A vertex of a polytope in $\mathbb R^{m}$ is
determined by $m$ linearly independent tight constraints, at most $|\mathcal P|$ of which are covering constraints,
so at most $|\mathcal P|$ coordinates of $y^*$ lie strictly between $0$ and $1$. Put $B=\{b:y^*_b>0\}$. Then
$\sum_{b\in B}\vp_p(b)\ge\sum_b\vp_p(b)y^*_b\ge R_p$ for every $p$, i.e.\ $\prod A\mid\prod B$, and
$|B|\le\#\{b:y^*_b=1\}+|\mathcal P|\le\tau^*+|\mathcal P|$.
\end{proof}

\begin{lemma}[few primes]\label{lem:fewprimes}
If $m<8n^3$ then $|\mathcal P|<16n$.
\end{lemma}
\begin{proof}
The $i$-th prime is at least $i+1$, so if $|\mathcal P|\ge16n$ then
$\prod_{p\in\mathcal P}p\ge(16n+1)!>(16n/e)^{16n}$, while $\prod_{p\in\mathcal P}p\le\prod A\le m^n<(8n^3)^n$.
Taking logarithms, $16\ln(16n/e)=28.3\ldots+16\ln n>\ln 8+3\ln n$ for every $n\ge1$, a contradiction.
\end{proof}

\begin{proof}[Proof of Theorem~\ref{thm:cond}]
If $a_n\ge8n^3$, Theorem~\ref{thm:long} gives $2n\le18n$ elements. Otherwise let $z$ be a maximiser in
Lemma~\ref{lem:dual} with $z_p\le1$, supported on $\mathcal P$. Since $w_z(a)\le2+(w_z(a)-2)^+$ and $A\subseteq[1,m]$,
\[
\sum_{a\in A}w_z(a)\le 2n+\sum_{k=1}^{m}\bigl(w_z(k)-2\bigr)^+\le 2n+\sum_{b\in I}\bigl(w_z(b)-1\bigr)^+
\]
by \eqref{eq:TH}, so $\tau^*\le2n$. Lemmas~\ref{lem:round} and~\ref{lem:fewprimes} give $|B|\le 2n+16n-1<18n$.
\end{proof}

\subsection*{Equivalent forms and a per-prime-power sufficient condition}
Subtracting $\sum w_z$ over both ranges, \eqref{eq:TH} is equivalent to
\begin{equation}\label{eq:TH2}
\sum_{b\in I}\min\bigl(w_z(b),1\bigr)\le\sum_{k=1}^{m}\min\bigl(w_z(k),2\bigr)+\Delta,\qquad
\Delta:=\sum_{b\in I}w_z(b)-\sum_{k=1}^{m}w_z(k)=\sum_{q}z_q\delta_q\ge0,
\end{equation}
where $q$ runs over prime powers, $z_{p^j}:=z_p$, and $\delta_q:=\#\{b\in I:q\mid b\}-\lfloor m/q\rfloor\in\{0,1\}$.
Writing $\min(w,1)=w\min(1,1/w)$ and $\min(w,2)=w\min(1,2/w)$ and expanding $w_z(n)=\sum_{q\mid n}z_q$, both sides of
\eqref{eq:TH2} become sums over prime powers, and \eqref{eq:TH} follows if for every prime power $q$ with $z_q>0$
\begin{equation}\label{eq:PQ}
\sum_{b\in I,\ q\mid b}\min\Bigl(1,\frac1{w_z(b)}\Bigr)\le\delta_q+\sum_{k\le m,\ q\mid k}\min\Bigl(1,\frac2{w_z(k)}\Bigr).
\end{equation}
Substituting $b=qi'$, $k=qi$ and $w_z(qi)=a+w_z(i)$ with $a:=w_z(q)$, the multiples of $q$ in $I$ correspond to an
interval $J$ of $L:=\lfloor m/q\rfloor$ or $L+1$ consecutive integers, and \eqref{eq:PQ} follows from
\begin{equation}\label{eq:PQprime}
\sum_{i\in J}f_a\bigl(w_z(i)\bigr)\le\sum_{i=1}^{L}g_a\bigl(w_z(i)\bigr),\qquad
f_a(u):=\min\Bigl(1,\frac1{a+u}\Bigr),\quad g_a(u):=\min\Bigl(1,\frac2{a+u}\Bigr),
\end{equation}
for intervals $J$ of length $L$ (the possible extra element of $J$ is paid by $\delta_q$, as $f_a\le1$).

\begin{proposition}\label{prop:PQcases}
Inequality \eqref{eq:PQprime} holds for every interval $J$ of length $L$ and every $a\ge0$ in each of the following
cases: (i) $z$ is supported on a single prime; (ii) $a\ge2$ and $aL\ge w_z(L!)$, in particular whenever
$a\ge\max\{2,\sum_{p\le L}z_p/(p-1)\}$.
\end{proposition}
\begin{proof}
(i) Let $z$ be supported on $p$ and $\phi(j):=f_a(z_pj)$, a nonincreasing function of $j$, with increments
$D_j:=\phi(j-1)-\phi(j)\ge0$. For any set $S$ of $L$ integers,
$\sum_{i\in S}\phi(\vp_p(i))=L\phi(0)-\sum_{j\ge1}D_j\#\{i\in S:p^j\mid i\}$. Since an interval of length $L$
contains at least $\lfloor L/p^j\rfloor=\#\{i\le L:p^j\mid i\}$ multiples of $p^j$, the sum over $J$ is at most the
sum over $[1,L]$, and $f_a\le g_a$ pointwise. (ii) For $a\ge2$ the minima are inactive; the left side is at most
$L/a$, and by the Cauchy--Schwarz inequality the right side is
$2\sum_{i\le L}(a+w_z(i))^{-1}\ge 2L^2/(aL+w_z(L!))\ge L/a$. Finally
$w_z(L!)=\sum_pz_p\vp_p(L!)\le L\sum_{p\le L}z_p/(p-1)$ by Legendre's formula.
\end{proof}

\subsection*{Three strengthenings that fail}
Each of the following statements would imply \eqref{eq:TH} or \eqref{eq:PQprime} by a short argument, and each is
false. The common mechanism is that an interval of length $L$ can contain more integers free of all primes $\le L$
than $\pi(L)+1$: the pattern of such integers is an admissible set, and by the Chinese remainder theorem every
admissible pattern occurs, so their maximal number is the Hensley--Richards function $\rho^*(L)$, which exceeds
$\pi(L)$ for $L\ge3159$ and by an unbounded amount as $L\to\infty$ \cite{HR74,GR98}. Explicit windows were built
by a centred sieve (residue class $0$ modulo every $p\le\sqrt L$, then the least populated class modulo each larger
prime) and the Chinese remainder theorem; the starting points $x$ and verification scripts are in the repository.

\begin{proposition}\label{prop:false}
\begin{enumerate}
\item[(a)] \emph{(weight-free transport)} There need not exist an injection $\varphi:[1,m]\to I$ such that for every
$k$ some prime $p\mid k$ satisfies $(k/p)\mid\varphi(k)$. For $m=20000$ there is an $x$ (with $4298$ digits) for which
$I$ contains $2270$ integers free of all primes $\le m/2$; every $k$ related to such an integer is $1$ or a prime, so
Hall's condition fails because $2270>\pi(m)+1=2263$. (Such an injection would give \eqref{eq:TH} at once, since
$w_z(\varphi(k))\ge w_z(k)-z_p\ge w_z(k)-1$.)
\item[(b)] \emph{(single-threshold level sets)} The inequality $\#\{i\le L:w_z(i)\ge2\tau\}\le\#\{i\in J:w_z(i)\ge\tau\}$
fails for $L=10^5$, $z_p=1$ for all $p\le L/2$, at $\tau=1$ (by $142$) and at $\tau=1/2$ (by $5275$), although
\eqref{eq:PQprime} holds on the same window with room to spare. (Integrating it over $\tau$ would give
\eqref{eq:PQprime}.)
\item[(c)] \emph{(Laplace form)} The inequality $s\sum_{i\in J}s^{w_z(i)}\le\sum_{i\le L}s^{w_z(i)}$ for all
$s\in(0,1]$, which would give \eqref{eq:PQprime} for all $a\ge2$ by integrating against $s^{a-1}\,ds$, fails for
$L=10^6$, $z_p=1$ for all $p\le L$, and $s\in[0.005,0.015]$: the window (with $x$ of $433636$ digits) contains
$80436$ integers free of all primes $\le L$, against $\pi(L)+1=78499$, and at $s=0.005$ the two sides are
$405.03$ and $398.77$.
\end{enumerate}
\end{proposition}

\begin{remark}[evidence and the remaining programme]
Conjecture~\ref{conj:TH} was tested on $2.1\times10^6$ random and adversarial instances with $m\le600$ (0/1,
fractional and $p^{-\theta}$ weights; $x=0$, random $x$, $x\equiv-1$ modulo products of weighted primes, windows centred
at highly divisible numbers) and, with nine weight choices, on the dense windows of Proposition~\ref{prop:false} of
lengths $2\cdot10^4$ and $10^5$, where the slack is at least $0.18m$; \eqref{eq:PQprime} was tested on $3.5\times10^5$
instances over all prime powers. Equality in \eqref{eq:TH} occurs only for $x=0$ and single primes. For 0/1 weights
$z=\mathbf 1_{\mathcal Q}$ the inequality is equivalent to $B_0-A_0\le\Delta+B_{\ge2}$, where $B_0$ and $A_0$ count the
integers in $[1,m]$ and in $I$ free of the primes of $\mathcal Q$, and $B_{\ge2}$ counts the $k\le m$ with at least two
prime factors from $\mathcal Q$ (with multiplicity); by inclusion--exclusion $B_0-A_0$ is an alternating sum of the
surpluses $\delta_d$ over squarefree products $d$ of primes of $\mathcal Q$. In the layer-cake form of
\eqref{eq:PQprime} the integrand is positive only where the count on $[1,L]$ is stuck at a low level, and its
positive part is bounded by the number of integers in a window free of the weighted primes, i.e.\ by an explicit large
sieve bound, while the negative part is a counting quantity on $[1,L]$. Making this quantitative for all weights is
the remaining step; the failures in Proposition~\ref{prop:false} show that it cannot be replaced by a weight-free
or single-threshold argument.
\end{remark}

\begin{thebibliography}{9}
\bibitem{ES59} P.~Erd\H{o}s and J.~Sur\'anyi, Megjegyz\'esek egy versenyfeladathoz (Remarks on a problem of a
mathematical competition), \emph{Mat.\ Lapok} 10 (1959), 39--48 (in Hungarian).
\bibitem{Er92} P.~Erd\H{o}s, Some of my forgotten problems in number theory, \emph{Hardy--Ramanujan J.} 15
(1992), 34--50, \S1.
\bibitem{EP708} T.~F.~Bloom, Erd\H{o}s Problem \#708, \url{https://www.erdosproblems.com/708}, accessed 2026-09-03.
\bibitem{vDLT26} W.~van Doorn, Y.~Li and Q.~Tang, Optimal bounds for an Erd\H{o}s problem on matching integers
to distinct multiples, arXiv:2603.28636 (2026).
\bibitem{HR74} D.~Hensley and I.~Richards, Primes in intervals, \emph{Acta Arith.} 25 (1974), 375--391.
\bibitem{GR98} D.~M.~Gordon and G.~Rodemich, Dense admissible sets, in: \emph{Algorithmic Number Theory (ANTS-III)}, Lecture Notes in Comput. Sci. 1423, Springer, 1998, 216--225.
\end{thebibliography}
\end{document}
